\section{Proofs for~\cref{sec:failure_uniform_opd}}
\label{app:proof_opd_failure}



We adopt the notation established in Section~\ref{sec:failure_uniform_opd}. For notational convenience in the derivations below, we write $p_t(\cdot) = \pi_\theta(\cdot \mid y_{<t})$ and $q_t(\cdot) = \pi_{\mathrm{teacher}}(\cdot \mid y_{<t})$ for the student and teacher next-token distributions conditioned on the student-generated prefix $y_{<t}$. The token-level OPD loss is $\ell_t = \log p_t(y_t) - \log q_t(y_t)$ with $y_t \sim p_t$. The step-level mismatch $\Delta_k$, the teacher-supported region $S_t^\epsilon$, and the overlap $\rho_t$ are as defined in the main text.

\subsection{Proof of Proposition~\ref{prop:discontinuous_drift}}

\begin{proof}
We analyze how the step-level mismatch $\Delta_k = \frac{1}{|\mathcal{I}_k|}\sum_{t \in \mathcal{I}_k} D_{\mathrm{KL}}(p_t \| q_t)$ evolves between consecutive steps in two settings.

\paragraph{Text-only case (gradual drift).}
In text-only generation, the prefix at position $t$ differs from position $t-1$ by exactly one student-sampled token. Define the per-token distributional shift as:
\begin{equation}
    \eta \triangleq \max_t \; \mathrm{TV}(p_t, p_{t-1}) = \max_t \; \frac{1}{2}\sum_{v \in \mathcal{V}}\left|p_t(v) - p_{t-1}(v)\right|,
\end{equation}
\ie the maximum total variation (TV) distance between consecutive output distributions. For well-trained transformers, $\eta$ is small because the self-attention mechanism distributes influence across many context positions, so appending a single token has limited impact on the output distribution. Since consecutive steps share a boundary of one token, the step-level mismatch changes by at most $O(\eta)$ per step transition.

\paragraph{TIR case (discontinuous drift).}
In TIR, the transition from step $k$ to step $k+1$ involves appending an entire tool observation $o_k = (o_k^1, \ldots, o_k^m)$ of length $m = |o_k|$ to the prefix:
\begin{equation}
    y_{<t_{k+1}^{\mathrm{start}}} = y_{<t_k^{\mathrm{end}}} \;\oplus\; o_k,
\end{equation}
where $\oplus$ denotes concatenation. Crucially, $o_k$ is produced by the external environment, not by either model. Let $\eta_{\mathrm{tool}}$ denote the average per-token distributional shift induced by observation tokens, measured by the TV distance between the model's output distribution before and after conditioning on each successive observation token.

\paragraph{Single-step bound.}
By the triangle inequality applied to the $m$ intermediate distributions between $p_{t_k^{\mathrm{end}}}$ and $p_{t_{k+1}^{\mathrm{start}}}$, the total variation shift satisfies:
\begin{equation}
    \mathrm{TV}\!\left(p_{t_{k+1}^{\mathrm{start}}},\; p_{t_k^{\mathrm{end}}}\right) \;\le\; m \cdot \eta_{\mathrm{tool}}.
\end{equation}
A similar bound holds for the teacher distribution. By Pinsker's inequality ($D_{\mathrm{KL}}(P\|Q) \ge 2\,\mathrm{TV}(P,Q)^2$), the change in step-level KL divergence is lower-bounded by the squared TV shift. Since tool observations typically span tens to hundreds of tokens ($m \gg 1$) and introduce content not generated by either model, we have $\eta_{\mathrm{tool}} > \eta$ in general, yielding:
\begin{equation}
    \Delta_{k+1} - \Delta_k \;=\; \Omega(m \cdot \eta_{\mathrm{tool}}) \;\ge\; O(\eta).
\end{equation}
Note that even a single erroneous tool call already introduces a substantially larger divergence jump than text-only drift ($\Omega(m \cdot \eta_{\mathrm{tool}})$ vs.\ $O(\eta)$), since the observation length $m \gg 1$ and the out-of-distribution content shifts conditioning more aggressively. Nevertheless, the teacher, having encountered some error patterns during training, can still provide partially useful supervision after an isolated failure. The critical issue arises from the \emph{cascading} nature of errors in weaker student models, as we analyze next.

\paragraph{Cascading errors and super-linear compounding.}
Weaker student models, precisely the targets of OPD, are prone to producing consecutive erroneous tool calls. A weaker student is more likely to generate an incorrect tool invocation at step $k$; conditioned on the resulting erroneous observation, its subsequent reasoning is further degraded, increasing the probability of another failure at step $k+1$. This creates a cascading failure pattern where each error compounds upon the previous ones.

The key insight is that while the teacher may reasonably handle a single isolated error in the context (having potentially seen similar error messages during training), the \emph{joint} occurrence of multiple consecutive failures creates a conditioning context that is combinatorially unlikely under the teacher's training distribution. If each individual error context has marginal probability $p_{\mathrm{err}}$ under the teacher's training distribution, the joint context of $j$ consecutive errors has probability at most $p_{\mathrm{err}}^j$, an exponential decay in familiarity. It is this accumulated, multi-error prefix that drives the teacher's conditional distribution far from calibration, making $\eta_{\mathrm{tool}}^{(i)}$ grow with $i$.

Formally, when the student makes consecutive erroneous tool calls across steps $k, k+1, \ldots, k+j-1$, the cumulative divergence satisfies:
\begin{equation}
    \Delta_{k+j} - \Delta_k \;=\; \Omega\!\left(\sum_{i=0}^{j-1} m_{k+i} \cdot \eta_{\mathrm{tool}}^{(k+i)}\right),
\end{equation}
where $m_{k+i}$ is the length of the $i$-th erroneous observation and $\eta_{\mathrm{tool}}^{(k+i)}$ is the corresponding per-token shift. Crucially, later errors induce progressively larger per-token shifts ($\eta_{\mathrm{tool}}^{(k+i+1)} \ge \eta_{\mathrm{tool}}^{(k+i)}$) because the prefix is already corrupted by prior errors, where the teacher has never been calibrated on such multi-error contexts, making its predictions increasingly unreliable. This yields super-linear growth of divergence with the number of consecutive tool failures, which is empirically confirmed in~\cref{motivation}(a) where the student-teacher divergence accelerates as errors accumulate.
\end{proof}

\subsection{Proof of Proposition~\ref{prop:variance_explosion}}

\begin{proof}
We prove the second-moment lower bound and then derive the SNR degradation.

\paragraph{Step 1: Decomposition of the second moment.}
The second moment of the OPD loss $\ell_t = \log p_t(y_t) - \log q_t(y_t)$ under $y_t \sim p_t$ decomposes over the vocabulary:
\begin{equation}
    \mathbb{E}_{y_t \sim p_t}[\ell_t^2]
    \;=\;
    \sum_{v \in \mathcal{V}} p_t(v) \left(\log \frac{p_t(v)}{q_t(v)}\right)^2.
\end{equation}

\paragraph{Step 2: Restricting to the low-overlap region.}
We partition the vocabulary into the teacher-supported region $S_t^\epsilon = \{v : q_t(v) \ge \epsilon\}$ and its complement $\bar{S}_t^\epsilon = \{v : q_t(v) < \epsilon\}$. By the overlap condition $\rho_t \le \rho$, the student places mass at least $1 - \rho$ on $\bar{S}_t^\epsilon$:
\begin{equation}
    \sum_{v \in \bar{S}_t^\epsilon} p_t(v) \;\ge\; 1 - \rho_t \;\ge\; 1 - \rho.
\end{equation}
Restricting the sum to $\bar{S}_t^\epsilon$ gives a lower bound:
\begin{equation}
    \mathbb{E}[\ell_t^2]
    \;\ge\;
    \sum_{v \in \bar{S}_t^\epsilon} p_t(v) \left(\log \frac{p_t(v)}{q_t(v)}\right)^2.
\end{equation}

\paragraph{Step 3: Applying Jensen's inequality.}
For any $v \in \bar{S}_t^\epsilon$, since $q_t(v) < \epsilon$, we have:
\begin{equation}
    \log \frac{p_t(v)}{q_t(v)} \;>\; \log p_t(v) + \log \frac{1}{\epsilon}.
\end{equation}
Define the conditional distribution $\tilde{p}(v) = p_t(v)/(1-\rho_t)$ over $\bar{S}_t^\epsilon$. By the convexity of $x \mapsto x^2$ and Jensen's inequality:
\begin{align}
    \mathbb{E}[\ell_t^2]
    &\;\ge\;
    (1-\rho_t) \sum_{v \in \bar{S}_t^\epsilon} \tilde{p}(v) \left(\log \frac{p_t(v)}{q_t(v)}\right)^2 \nonumber \\
    &\;\ge\;
    (1-\rho_t) \left(
    \sum_{v \in \bar{S}_t^\epsilon} \tilde{p}(v) \log \frac{p_t(v)}{q_t(v)}
    \right)^2 \nonumber \\
    &\;\ge\;
    (1-\rho_t) \left(
    \log \frac{1}{\epsilon} + \sum_{v \in \bar{S}_t^\epsilon} \tilde{p}(v) \log p_t(v)
    \right)^2.
\end{align}

\paragraph{Step 4: Bounding the entropy term.}
The term $\sum_{v \in \bar{S}_t^\epsilon} \tilde{p}(v) \log p_t(v) = -H_{\bar{S}}$, where $H_{\bar{S}}$ is the conditional entropy of the student restricted to $\bar{S}_t^\epsilon$, satisfying $0 \le H_{\bar{S}} \le \log|\mathcal{V}|$. For sufficiently small $\epsilon$ satisfying $\epsilon < 1/|\mathcal{V}|$ (a natural condition since the teacher threshold should be below the uniform baseline), we have $\log(1/\epsilon) > \log|\mathcal{V}| \ge H_{\bar{S}}$, so the $\log(1/\epsilon)$ term dominates:
\begin{equation}
    \mathbb{E}[\ell_t^2]
    \;\ge\;
    (1-\rho)\left(\log \frac{1}{\epsilon} - H_{\bar{S}}\right)^2
    \;\ge\;
    (1-\rho)\,\log^2\!\frac{1}{\epsilon} \cdot \left(1 - \frac{\log|\mathcal{V}|}{\log(1/\epsilon)}\right)^2,
\end{equation}
which grows as $\Omega(\log^2(1/\epsilon))$ when $\rho_t \to 0$.

\paragraph{Step 5: SNR degradation.}
Let $g_t = \ell_t \nabla_\theta \log \pi_\theta(y_t \mid y_{<t})$ be the token-level OPD gradient estimator. The expected gradient is $\mathbb{E}[g_t] = \nabla_\theta D_{\mathrm{KL}}(p_t \| q_t)$, which is bounded for any fixed distribution pair with finite KL divergence; let $C = \|\mathbb{E}[g_t]\|$.

Define $c_{\min} \triangleq \min_v \|\nabla_\theta \log \pi_\theta(v \mid y_{<t})\|^2 > 0$, which is positive for any non-degenerate parameterization. The variance of the gradient estimator satisfies:
\begin{equation}
    \mathrm{Var}[\|g_t\|]
    \;\ge\;
    \mathbb{E}[\ell_t^2] \cdot c_{\min} - C^2.
\end{equation}
Combining with the second-moment lower bound from Step 4:
\begin{equation}
    \mathrm{SNR}(g_t)
    \;=\;
    \frac{\|\mathbb{E}[g_t]\|}{\sqrt{\mathrm{Var}[g_t]}}
    \;\le\;
    \frac{C}{\sqrt{(1-\rho)\log^2(1/\epsilon) \cdot c_{\min} - C^2}}.
\end{equation}
As $\rho \to 0$, the denominator grows as $O(\log(1/\epsilon))$ while the numerator $C$ remains constant, giving $\mathrm{SNR}(g_t) \to 0$.

\paragraph{Interpretation.} When the student drifts far from the teacher-supported region ($\rho_t \approx 0$), most sampled tokens fall in $\bar{S}_t^\epsilon$ where the teacher assigns near-zero probability. The OPD loss for these tokens has large magnitude (due to $-\log q_t(y_t) > \log(1/\epsilon)$) but is uninformative. It reflects the teacher's inability to model these out-of-distribution states rather than a meaningful learning signal. Meanwhile, informative signals from $S_t^\epsilon$ (where teacher guidance is calibrated) are sampled only with probability $\rho_t$, making them increasingly rare. The gradient is thus dominated by high-variance, low-information contributions, rendering optimization unstable.
\end{proof}


\subsection{Discussion: Failure Cascade and Design Implications}
\label{app:discussion_objective}

The two propositions together characterize a failure cascade specific to TIR:
\begin{enumerate}
    \item \textbf{Initial perturbation}: An erroneous tool call returns a corrupted observation (\eg a runtime error, incorrect output, or timeout message). This already introduces a divergence jump substantially larger than text-only drift ($\Omega(m \cdot \eta_{\mathrm{tool}})$ vs.\ $O(\eta)$), though the teacher, having encountered some error patterns during pretraining, can still provide partially useful supervision after an isolated failure.
    
    \item \textbf{Cascading accumulation}: Weaker student models, precisely the targets of OPD, are prone to making consecutive errors. Each subsequent erroneous tool call further corrupts the prefix, and the \emph{joint} pattern of multiple consecutive failures becomes exponentially unlikely under the teacher's training distribution ($\sim p_{\mathrm{err}}^j$ for $j$ consecutive errors). It is this accumulated multi-error context, rather than any single error, that drives the teacher's conditional distribution far from calibration. The divergence thus compounds super-linearly (Proposition~\ref{prop:discontinuous_drift}), as empirically demonstrated by the accelerating divergence curve in~\cref{motivation}(a).
    
    \item \textbf{Supervision breakdown}: In the resulting low-overlap states ($\rho_t \approx 0$) caused by accumulated consecutive errors, the OPD gradient estimator suffers variance explosion and SNR degradation (Proposition~\ref{prop:variance_explosion}). Updates become dominated by uninformative, high-magnitude contributions from tokens where the teacher provides no meaningful guidance.~\cref{motivation}(b) confirms this empirically: the teacher's conditional entropy becomes both elevated and highly variable in steps following accumulated tool errors.
    
    \item \textbf{Amplification by uniform aggregation}: OPD sums token-level losses across all steps with equal weight, treating well-aligned early steps (high $\rho_t$, reliable teacher guidance) identically to corrupted post-error steps (low $\rho_t$, unreliable guidance). When cumulative tool errors cause a substantial portion of later steps to have low overlap, the aggregate gradient is systematically biased by these high-variance contributions.
\end{enumerate}

This analysis implies that a stabilizing objective should modulate distillation strength at the step level: preserving full-strength dense supervision in well-aligned steps while attenuating the signal when the estimated student-teacher divergence indicates teacher miscalibration. This is the design principle underlying our adaptive step-wise reweighting mechanism.















\subsection{Variance Reduction under Step-wise Reweighting}
\label{app:variance_reduction}

We now show that the step-wise reweighting mechanism in Eq.~\eqref{eq:adaptive_weight} provably addresses the gradient SNR degradation identified in Proposition~\ref{prop:variance_explosion}.

\begin{proposition}[Bounded variance under adaptive reweighting]
\label{prop:variance_reduction}
Under the step-wise weighted OPD objective (Eq.~\ref{eq:step_opd}), when the divergence increases monotonically across steps ($d_1 \le d_2 \le \cdots \le d_k$), the weighted gradient contribution from step $k$ satisfies:
\begin{equation}
    \mathbb{E}[w_k^2 \cdot \ell_t^2] \;\le\; \frac{(d_1 + \epsilon)^2}{(d_k + \epsilon)^2} \cdot \mathbb{E}[\ell_t^2],
\end{equation}
where $\mathbb{E}[\ell_t^2]$ is the unweighted second moment from Proposition~\ref{prop:variance_explosion}. Consequently, even when $\rho_t \to 0$ causes $\mathbb{E}[\ell_t^2] \to \infty$, the weighted contribution remains bounded whenever $d_k$ grows proportionally to the divergence.
\end{proposition}

\begin{proof}
By the definition of $w_k$ in Eq.~\eqref{eq:adaptive_weight}:
\begin{equation}
    w_k = \min\!\left(\prod_{u=1}^{k-1} \frac{d_u + \epsilon}{d_{u+1} + \epsilon},\; 1+\delta\right).
\end{equation}
Under monotonically increasing divergence ($d_1 \le d_2 \le \cdots \le d_k$), each ratio $\frac{d_u + \epsilon}{d_{u+1} + \epsilon} \le 1$, and the product telescopes:
\begin{equation}
    w_k \;\le\; \prod_{u=1}^{k-1} \frac{d_u + \epsilon}{d_{u+1} + \epsilon} \;=\; \frac{d_1 + \epsilon}{d_k + \epsilon}.
\end{equation}
The weighted second moment of the OPD loss at token $t \in \mathcal{I}_k$ is therefore:
\begin{equation}
    \mathbb{E}[w_k^2 \cdot \ell_t^2] \;\le\; \left(\frac{d_1 + \epsilon}{d_k + \epsilon}\right)^2 \cdot \mathbb{E}[\ell_t^2].
\end{equation}

From Proposition~\ref{prop:variance_explosion}, when the student drifts into low-overlap states, $\mathbb{E}[\ell_t^2] \ge (1-\rho)\log^2(1/\epsilon_{\mathrm{teacher}})$. However, such drift also implies that $d_k$ increases (since $d_k$ is monotonically related to $\Delta_k$, as shown in Appendix~\ref{app:dk_proxy}). Specifically, when $\rho_t \to 0$, the per-token log-probability gaps grow, causing $d_k \gg d_1$. The weighted contribution thus satisfies:
\begin{equation}
    \mathbb{E}[w_k^2 \cdot \ell_t^2] \;\le\; \frac{(d_1 + \epsilon)^2}{(d_k + \epsilon)^2} \cdot (1-\rho)\log^2(1/\epsilon_{\mathrm{teacher}}).
\end{equation}

The key insight is that the numerator $(d_1 + \epsilon)^2$ is bounded (determined by the initial step's divergence), while the denominator $(d_k + \epsilon)^2$ grows with the accumulated drift. This creates an automatic variance suppression: the more the student diverges (larger $d_k$, larger $\mathbb{E}[\ell_t^2]$), the more aggressively the weight $w_k$ attenuates the contribution, preventing the variance explosion that afflicts OPD.

\paragraph{SNR recovery.}
For the weighted gradient estimator $\tilde{g}_t = w_k \cdot \ell_t \cdot \nabla_\theta \log \pi_\theta(y_t \mid y_{<t})$, the signal (expected gradient from well-aligned early steps where $w_k \approx 1$) remains intact, while the noise from high-divergence later steps is suppressed by the factor $(d_1 + \epsilon)^2/(d_k + \epsilon)^2$. This restores a positive SNR for the aggregate gradient, in contrast to OPD where the SNR degrades to zero (Proposition~\ref{prop:variance_explosion}).
\end{proof}

\paragraph{Stability under recovery.}
The above analysis addresses the erroneous case where divergence increases monotonically. In the recovery case ($d_{u+1} < d_u$ for some $u$), the weight $w_k$ may exceed 1, potentially amplifying the gradient contribution. However, the hard upper bound $w_k \le 1 + \delta$ in Eq.~\eqref{eq:adaptive_weight} ensures that the weighted second moment is always bounded by $(1+\delta)^2 \cdot \mathbb{E}[\ell_t^2]$. Since recovery steps by definition have decreasing $d_k$ (implying the student is returning to the teacher-supported region, \ie $\rho_t$ is increasing), the unweighted $\mathbb{E}[\ell_t^2]$ itself is decreasing in these steps. The combination of bounded amplification and decreasing base variance ensures optimization stability.

\paragraph{Interpretation.}
This result formalizes the intuition behind our method: OPD fails because it applies uniform weight to all steps, allowing high-variance contributions from corrupted post-error steps to dominate the gradient. Our step-wise reweighting automatically detects these corrupted steps (via increasing $d_k$) and attenuates their influence proportionally, ensuring that the gradient signal remains dominated by informative contributions from well-aligned steps. The upper bound $1+\delta$ further ensures that the recovery mechanism does not over-amplify any single step, maintaining optimization stability throughout training.

















\subsection{Justification of $d_k$ as a Proxy for $\Delta_k$}
\label{app:dk_proxy}

Computing the full KL divergence $\Delta_k$ requires evaluating the teacher's output distribution over the entire vocabulary at each token position, which introduces substantial overhead. In contrast, $d_k$ (Eq.~\ref{eq:step_div}) only requires the teacher's log-probability on the student-sampled token $y_t$, a quantity already computed in the standard OPD forward pass.

\paragraph{Monotonic consistency.}
By Jensen's inequality:
\begin{equation}
    d_k = \frac{1}{|\mathcal{I}_k|} \sum_{t \in \mathcal{I}_k} \left| \log \frac{\pi_\theta(y_t \mid y_{<t})}{\pi_{\mathrm{teacher}}(y_t \mid y_{<t})} \right|
    \;\ge\;
    \left| \frac{1}{|\mathcal{I}_k|} \sum_{t \in \mathcal{I}_k} \log \frac{\pi_\theta(y_t \mid y_{<t})}{\pi_{\mathrm{teacher}}(y_t \mid y_{<t})} \right|,
\end{equation}
where the right-hand side is the absolute value of a single-sample Monte Carlo estimate of $\Delta_k$. When the student drifts further from the teacher (increasing $\Delta_k$), the per-token log-probability gaps increase in expectation, yielding a larger $d_k$. This monotonic relationship ensures that the ratios $d_u / d_{u+1}$ used in Eq.~\eqref{eq:adaptive_weight} correctly reflect the relative trend of divergence across steps.

\paragraph{Sufficiency of ordering.}
The weight formula in Eq.~\eqref{eq:adaptive_weight} is a product of ratios $\frac{d_u + \epsilon}{d_{u+1} + \epsilon}$, depending only on the relative magnitudes between consecutive steps. Any monotone transformation of $\Delta_k$ preserves these ratios' direction (above or below one), producing equivalent attenuation and recovery behavior. Since $d_k$ maintains the ordering of $\Delta_k$, it is a sufficient statistic for our reweighting mechanism.
